\scanpage{117}
(2) \deleted{\unsure{Posons} $u_a$}
\inserted{Soit $v^a_{\unsure{\xi|\zeta}}$ le composé}
\diagram{\begin{tikzcd}[column sep=large]
\unsure{\xi|\zeta} \arrow[r,"u_{a,\unsure{\xi|\zeta}}"] & a\unsure{\xi|\zeta} \arrow[r,"u_{a,a\unsure{\xi|\zeta}}"] & a^2\unsure{\xi|\zeta}=\unsure{\xi|\zeta}
\end{tikzcd}}
alors les $v^a_{\unsure{\xi|\zeta}}$ commutent ($\unsure{\xi|\zeta}$ fixé, $a\in\widehat{A}$ \unsure{variab.}).

\margin{\uline{\unsure{NB}}}
Le groupoïde libre sur les \deleted{\unsure{objets de}} \inserted{$\varepsilon_G$ comme} ens. d'objets, et avec les générateurs $u_{a,\unsure{\xi|\zeta}}$ et les relations (1) (2), \deleted{est} s'identifie bien à $S_0\mathcal{T}_G$ (à isom. can. près).

\unsure{Prenant} $\mathcal{T}_G$, au \deleted{\unsure{point}} lieu de $S_0\mathcal{T}_G$, on trouve une construction \deleted{\unsure{analogue}} \inserted{\unsure{parallèle}} \deleted{en} termes de \deleted{$\varepsilon_{G,A}$ \unsure{(quotient de} $\varepsilon$}
\[
\varepsilon_G=\prod_{a\in A}\varepsilon_a\subset P_G=\prod_{a\in A}T_a
\]
qui est l'ens des objets d'un groupoïde \inserted{\unsure{conn.}} ${}_0\mathcal{T}_G$, \deleted{décrit} de groupe fondamental abélien $\mathbb{Z}^A$, décrit mutatis mutandis, comme $S_0\mathcal{T}_G$ à partir des $\unsure{S\varepsilon_G|S_{\varepsilon_G}}$. On a alors
\diagram{\begin{tikzcd}[column sep=large,row sep=large]
S_0\mathcal{T}_G \arrow[r] \arrow[d] & S\mathcal{T}_{G^{\unsure{\natural|\sharp}}} \arrow[d] \\
{}_0\mathcal{T}_G \arrow[r] & \mathcal{T}_{G^{\unsure{\natural|\sharp}}}
\end{tikzcd}}
\mspage{116}
